\documentclass[]{article}
\usepackage{lmodern}
\usepackage{amssymb,amsmath}
\usepackage{ifxetex,ifluatex}
\usepackage{fixltx2e} % provides \textsubscript
\ifnum 0\ifxetex 1\fi\ifluatex 1\fi=0 % if pdftex
  \usepackage[T1]{fontenc}
  \usepackage[utf8]{inputenc}
\else % if luatex or xelatex
  \ifxetex
    \usepackage{mathspec}
  \else
    \usepackage{fontspec}
  \fi
  \defaultfontfeatures{Ligatures=TeX,Scale=MatchLowercase}
\fi
% use upquote if available, for straight quotes in verbatim environments
\IfFileExists{upquote.sty}{\usepackage{upquote}}{}
% use microtype if available
\IfFileExists{microtype.sty}{%
\usepackage{microtype}
\UseMicrotypeSet[protrusion]{basicmath} % disable protrusion for tt fonts
}{}
\usepackage[left=1.1in, right=1.1in, top=1.1in, bottom=1.1in]{geometry}
\usepackage{hyperref}
\hypersetup{unicode=true,
            pdftitle={Stat 435 Lab Notes 1},
            pdfauthor={Xiongzhi Chen; Washington State University},
            pdfborder={0 0 0},
            breaklinks=true}
\urlstyle{same}  % don't use monospace font for urls
\usepackage{color}
\usepackage{fancyvrb}
\newcommand{\VerbBar}{|}
\newcommand{\VERB}{\Verb[commandchars=\\\{\}]}
\DefineVerbatimEnvironment{Highlighting}{Verbatim}{commandchars=\\\{\}}
% Add ',fontsize=\small' for more characters per line
\usepackage{framed}
\definecolor{shadecolor}{RGB}{248,248,248}
\newenvironment{Shaded}{\begin{snugshade}}{\end{snugshade}}
\newcommand{\KeywordTok}[1]{\textcolor[rgb]{0.13,0.29,0.53}{\textbf{{#1}}}}
\newcommand{\DataTypeTok}[1]{\textcolor[rgb]{0.13,0.29,0.53}{{#1}}}
\newcommand{\DecValTok}[1]{\textcolor[rgb]{0.00,0.00,0.81}{{#1}}}
\newcommand{\BaseNTok}[1]{\textcolor[rgb]{0.00,0.00,0.81}{{#1}}}
\newcommand{\FloatTok}[1]{\textcolor[rgb]{0.00,0.00,0.81}{{#1}}}
\newcommand{\ConstantTok}[1]{\textcolor[rgb]{0.00,0.00,0.00}{{#1}}}
\newcommand{\CharTok}[1]{\textcolor[rgb]{0.31,0.60,0.02}{{#1}}}
\newcommand{\SpecialCharTok}[1]{\textcolor[rgb]{0.00,0.00,0.00}{{#1}}}
\newcommand{\StringTok}[1]{\textcolor[rgb]{0.31,0.60,0.02}{{#1}}}
\newcommand{\VerbatimStringTok}[1]{\textcolor[rgb]{0.31,0.60,0.02}{{#1}}}
\newcommand{\SpecialStringTok}[1]{\textcolor[rgb]{0.31,0.60,0.02}{{#1}}}
\newcommand{\ImportTok}[1]{{#1}}
\newcommand{\CommentTok}[1]{\textcolor[rgb]{0.56,0.35,0.01}{\textit{{#1}}}}
\newcommand{\DocumentationTok}[1]{\textcolor[rgb]{0.56,0.35,0.01}{\textbf{\textit{{#1}}}}}
\newcommand{\AnnotationTok}[1]{\textcolor[rgb]{0.56,0.35,0.01}{\textbf{\textit{{#1}}}}}
\newcommand{\CommentVarTok}[1]{\textcolor[rgb]{0.56,0.35,0.01}{\textbf{\textit{{#1}}}}}
\newcommand{\OtherTok}[1]{\textcolor[rgb]{0.56,0.35,0.01}{{#1}}}
\newcommand{\FunctionTok}[1]{\textcolor[rgb]{0.00,0.00,0.00}{{#1}}}
\newcommand{\VariableTok}[1]{\textcolor[rgb]{0.00,0.00,0.00}{{#1}}}
\newcommand{\ControlFlowTok}[1]{\textcolor[rgb]{0.13,0.29,0.53}{\textbf{{#1}}}}
\newcommand{\OperatorTok}[1]{\textcolor[rgb]{0.81,0.36,0.00}{\textbf{{#1}}}}
\newcommand{\BuiltInTok}[1]{{#1}}
\newcommand{\ExtensionTok}[1]{{#1}}
\newcommand{\PreprocessorTok}[1]{\textcolor[rgb]{0.56,0.35,0.01}{\textit{{#1}}}}
\newcommand{\AttributeTok}[1]{\textcolor[rgb]{0.77,0.63,0.00}{{#1}}}
\newcommand{\RegionMarkerTok}[1]{{#1}}
\newcommand{\InformationTok}[1]{\textcolor[rgb]{0.56,0.35,0.01}{\textbf{\textit{{#1}}}}}
\newcommand{\WarningTok}[1]{\textcolor[rgb]{0.56,0.35,0.01}{\textbf{\textit{{#1}}}}}
\newcommand{\AlertTok}[1]{\textcolor[rgb]{0.94,0.16,0.16}{{#1}}}
\newcommand{\ErrorTok}[1]{\textcolor[rgb]{0.64,0.00,0.00}{\textbf{{#1}}}}
\newcommand{\NormalTok}[1]{{#1}}
\usepackage{graphicx,grffile}
\makeatletter
\def\maxwidth{\ifdim\Gin@nat@width>\linewidth\linewidth\else\Gin@nat@width\fi}
\def\maxheight{\ifdim\Gin@nat@height>\textheight\textheight\else\Gin@nat@height\fi}
\makeatother
% Scale images if necessary, so that they will not overflow the page
% margins by default, and it is still possible to overwrite the defaults
% using explicit options in \includegraphics[width, height, ...]{}
\setkeys{Gin}{width=\maxwidth,height=\maxheight,keepaspectratio}
\IfFileExists{parskip.sty}{%
\usepackage{parskip}
}{% else
\setlength{\parindent}{0pt}
\setlength{\parskip}{6pt plus 2pt minus 1pt}
}
\setlength{\emergencystretch}{3em}  % prevent overfull lines
\providecommand{\tightlist}{%
  \setlength{\itemsep}{0pt}\setlength{\parskip}{0pt}}
\setcounter{secnumdepth}{0}
% Redefines (sub)paragraphs to behave more like sections
\ifx\paragraph\undefined\else
\let\oldparagraph\paragraph
\renewcommand{\paragraph}[1]{\oldparagraph{#1}\mbox{}}
\fi
\ifx\subparagraph\undefined\else
\let\oldsubparagraph\subparagraph
\renewcommand{\subparagraph}[1]{\oldsubparagraph{#1}\mbox{}}
\fi

%%% Use protect on footnotes to avoid problems with footnotes in titles
\let\rmarkdownfootnote\footnote%
\def\footnote{\protect\rmarkdownfootnote}

%%% Change title format to be more compact
\usepackage{titling}

% Create subtitle command for use in maketitle
\newcommand{\subtitle}[1]{
  \posttitle{
    \begin{center}\large#1\end{center}
    }
}

\setlength{\droptitle}{-2em}

  \title{Stat 435 Lab Notes 1}
    \pretitle{\vspace{\droptitle}\centering\huge}
  \posttitle{\par}
    \author{Xiongzhi Chen \\ Washington State University}
    \preauthor{\centering\large\emph}
  \postauthor{\par}
    \date{}
    \predate{}\postdate{}
  
\usepackage{bbm}
\usepackage{amssymb}
\usepackage{amsmath}
\usepackage{graphicx,float}

\begin{document}
\maketitle

{
\setcounter{tocdepth}{2}
\tableofcontents
}
\section{}\label{section}

\section{Basic objects and their
operations}\label{basic-objects-and-their-operations}

\subsection{Important objects}\label{important-objects}

Some important objects in R

\begin{itemize}
\tightlist
\item
  \texttt{vector}
\item
  \texttt{matrix} and \texttt{data.frame}
\item
  \texttt{character} and \texttt{factor}
\end{itemize}

In terms of data analysis, \texttt{data.frame} is perhaps the most
important R object.

\subsection{\texorpdfstring{\texttt{vector}}{vector}}\label{vector}

Vectors can be created by concatenation (\texttt{c}), sequencing
(\texttt{:} or \texttt{seq}), or replication (\texttt{rep}):

\begin{Shaded}
\begin{Highlighting}[]
\NormalTok{>}\StringTok{ }\CommentTok{# create a vector via concatenation}
\ErrorTok{>}\StringTok{ }\NormalTok{x =}\StringTok{ }\KeywordTok{c}\NormalTok{(}\DecValTok{1}\NormalTok{,}\DecValTok{5}\NormalTok{,}\FloatTok{0.8}\NormalTok{,}\DecValTok{10}\NormalTok{)}
\NormalTok{>}\StringTok{ }\NormalTok{x}
\NormalTok{[}\DecValTok{1}\NormalTok{]  }\FloatTok{1.0}  \FloatTok{5.0}  \FloatTok{0.8} \FloatTok{10.0}
\NormalTok{>}\StringTok{ }\CommentTok{# create a vector via seq}
\ErrorTok{>}\StringTok{ }\NormalTok{y=}\DecValTok{1}\NormalTok{:}\DecValTok{3}
\NormalTok{>}\StringTok{ }\NormalTok{y}
\NormalTok{[}\DecValTok{1}\NormalTok{] }\DecValTok{1} \DecValTok{2} \DecValTok{3}
\NormalTok{>}\StringTok{ }\NormalTok{z =}\StringTok{ }\KeywordTok{seq}\NormalTok{(}\DataTypeTok{from=}\DecValTok{1}\NormalTok{,}\DataTypeTok{to=}\DecValTok{10}\NormalTok{,}\DataTypeTok{by=}\DecValTok{2}\NormalTok{)}
\NormalTok{>}\StringTok{ }\NormalTok{z}
\NormalTok{[}\DecValTok{1}\NormalTok{] }\DecValTok{1} \DecValTok{3} \DecValTok{5} \DecValTok{7} \DecValTok{9}
\NormalTok{>}\StringTok{ }\CommentTok{# create a vector via rep}
\ErrorTok{>}\StringTok{ }\NormalTok{w =}\StringTok{ }\KeywordTok{rep}\NormalTok{(}\DecValTok{1}\NormalTok{,}\DataTypeTok{times=}\DecValTok{5}\NormalTok{)}
\NormalTok{>}\StringTok{ }\NormalTok{w}
\NormalTok{[}\DecValTok{1}\NormalTok{] }\DecValTok{1} \DecValTok{1} \DecValTok{1} \DecValTok{1} \DecValTok{1}
\NormalTok{>}\StringTok{ }\NormalTok{u =}\StringTok{ }\KeywordTok{rep}\NormalTok{(}\KeywordTok{c}\NormalTok{(}\DecValTok{1}\NormalTok{,}\DecValTok{2}\NormalTok{),}\DataTypeTok{each=}\DecValTok{3}\NormalTok{)}
\NormalTok{>}\StringTok{ }\NormalTok{u}
\NormalTok{[}\DecValTok{1}\NormalTok{] }\DecValTok{1} \DecValTok{1} \DecValTok{1} \DecValTok{2} \DecValTok{2} \DecValTok{2}
\end{Highlighting}
\end{Shaded}

\subsection{\texorpdfstring{\texttt{vector}}{vector}}\label{vector-1}

Entries of a vector can be accessed by specifying their indices as a
vector:

\begin{Shaded}
\begin{Highlighting}[]
\NormalTok{>}\StringTok{ }\NormalTok{x =}\StringTok{ }\KeywordTok{c}\NormalTok{(}\DecValTok{1}\NormalTok{,}\DecValTok{5}\NormalTok{,}\FloatTok{0.8}\NormalTok{,}\DecValTok{10}\NormalTok{)}
\NormalTok{>}\StringTok{ }\NormalTok{x}
\NormalTok{[}\DecValTok{1}\NormalTok{]  }\FloatTok{1.0}  \FloatTok{5.0}  \FloatTok{0.8} \FloatTok{10.0}
\NormalTok{>}\StringTok{ }\CommentTok{# obtain number of entries of a vector}
\ErrorTok{>}\StringTok{ }\KeywordTok{length}\NormalTok{(x)}
\NormalTok{[}\DecValTok{1}\NormalTok{] }\DecValTok{4}
\NormalTok{>}\StringTok{ }\CommentTok{# access entries of a vector}
\ErrorTok{>}\StringTok{ }\NormalTok{x[}\DecValTok{1}\NormalTok{]}
\NormalTok{[}\DecValTok{1}\NormalTok{] }\DecValTok{1}
\NormalTok{>}\StringTok{ }\NormalTok{x[}\DecValTok{1}\NormalTok{:}\DecValTok{3}\NormalTok{]}
\NormalTok{[}\DecValTok{1}\NormalTok{] }\FloatTok{1.0} \FloatTok{5.0} \FloatTok{0.8}
\NormalTok{>}\StringTok{ }\CommentTok{# remove some entries}
\ErrorTok{>}\StringTok{ }\NormalTok{y=x[-}\KeywordTok{c}\NormalTok{(}\DecValTok{1}\NormalTok{,}\DecValTok{3}\NormalTok{)]}
\NormalTok{>}\StringTok{ }\NormalTok{y}
\NormalTok{[}\DecValTok{1}\NormalTok{]  }\DecValTok{5} \DecValTok{10}
\end{Highlighting}
\end{Shaded}

\subsection{\texorpdfstring{\texttt{matrix} and
\texttt{data.frame}}{matrix and data.frame}}\label{matrix-and-data.frame}

Both \texttt{data.frame} and \texttt{matrix} have a grid layout in rows
and columns. However, the object types for all columns of a
\texttt{matrix} have to be the same.

\begin{itemize}
\tightlist
\item
  A matrix can be created by the command \texttt{matrix}, or by
  combining vectors as rows (\texttt{rbind}) or as columns
  (\texttt{cbind})
\item
  A \texttt{matrix} can be coerced into a data frame via
  \texttt{as.data.frame}
\item
  A data frame can be created by the command \texttt{data.frame} or
  automatically when importing data
\end{itemize}

\subsection{\texorpdfstring{\texttt{matrix} and
\texttt{data.frame}}{matrix and data.frame}}\label{matrix-and-data.frame-1}

\begin{Shaded}
\begin{Highlighting}[]
\NormalTok{>}\StringTok{ }\CommentTok{# create a matrix}
\ErrorTok{>}\StringTok{ }\NormalTok{x =}\StringTok{ }\KeywordTok{matrix}\NormalTok{(}\DecValTok{1}\NormalTok{:}\DecValTok{10}\NormalTok{,}\DataTypeTok{nrow=}\DecValTok{2}\NormalTok{,}\DataTypeTok{ncol=}\DecValTok{5}\NormalTok{)}
\NormalTok{>}\StringTok{ }\NormalTok{x}
     \NormalTok{[,}\DecValTok{1}\NormalTok{] [,}\DecValTok{2}\NormalTok{] [,}\DecValTok{3}\NormalTok{] [,}\DecValTok{4}\NormalTok{] [,}\DecValTok{5}\NormalTok{]}
\NormalTok{[}\DecValTok{1}\NormalTok{,]    }\DecValTok{1}    \DecValTok{3}    \DecValTok{5}    \DecValTok{7}    \DecValTok{9}
\NormalTok{[}\DecValTok{2}\NormalTok{,]    }\DecValTok{2}    \DecValTok{4}    \DecValTok{6}    \DecValTok{8}   \DecValTok{10}
\NormalTok{>}\StringTok{ }\KeywordTok{dim}\NormalTok{(x)}
\NormalTok{[}\DecValTok{1}\NormalTok{] }\DecValTok{2} \DecValTok{5}
\NormalTok{>}\StringTok{ }\CommentTok{# convert x into a data.frame}
\ErrorTok{>}\StringTok{ }\NormalTok{y =}\StringTok{ }\KeywordTok{as.data.frame}\NormalTok{(x)}
\NormalTok{>}\StringTok{ }\NormalTok{y}
  \NormalTok{V1 V2 V3 V4 V5}
\DecValTok{1}  \DecValTok{1}  \DecValTok{3}  \DecValTok{5}  \DecValTok{7}  \DecValTok{9}
\DecValTok{2}  \DecValTok{2}  \DecValTok{4}  \DecValTok{6}  \DecValTok{8} \DecValTok{10}
\NormalTok{>}\StringTok{ }\KeywordTok{dim}\NormalTok{(y)}
\NormalTok{[}\DecValTok{1}\NormalTok{] }\DecValTok{2} \DecValTok{5}
\NormalTok{>}\StringTok{ }\CommentTok{# indexing for matrix or data.frame}
\ErrorTok{>}\StringTok{ }\NormalTok{x[}\KeywordTok{c}\NormalTok{(}\DecValTok{1}\NormalTok{,}\DecValTok{2}\NormalTok{),}\KeywordTok{c}\NormalTok{(}\DecValTok{1}\NormalTok{,}\DecValTok{5}\NormalTok{)]}
     \NormalTok{[,}\DecValTok{1}\NormalTok{] [,}\DecValTok{2}\NormalTok{]}
\NormalTok{[}\DecValTok{1}\NormalTok{,]    }\DecValTok{1}    \DecValTok{9}
\NormalTok{[}\DecValTok{2}\NormalTok{,]    }\DecValTok{2}   \DecValTok{10}
\NormalTok{>}\StringTok{ }\NormalTok{y[}\KeywordTok{c}\NormalTok{(}\DecValTok{1}\NormalTok{,}\DecValTok{2}\NormalTok{),}\KeywordTok{c}\NormalTok{(}\DecValTok{1}\NormalTok{,}\DecValTok{5}\NormalTok{)]}
  \NormalTok{V1 V5}
\DecValTok{1}  \DecValTok{1}  \DecValTok{9}
\DecValTok{2}  \DecValTok{2} \DecValTok{10}
\NormalTok{>}\StringTok{ }\CommentTok{# access variable in data frame via $}
\ErrorTok{>}\StringTok{ }\NormalTok{y$V1}
\NormalTok{[}\DecValTok{1}\NormalTok{] }\DecValTok{1} \DecValTok{2}
\end{Highlighting}
\end{Shaded}

\subsection{\texorpdfstring{\texttt{matrix} and
\texttt{data.frame}}{matrix and data.frame}}\label{matrix-and-data.frame-2}

\begin{Shaded}
\begin{Highlighting}[]
\NormalTok{>}\StringTok{ }\NormalTok{x =}\StringTok{ }\DecValTok{1}\NormalTok{:}\DecValTok{3}\NormalTok{; y =}\StringTok{ }\DecValTok{2}\NormalTok{:}\DecValTok{4}
\NormalTok{>}\StringTok{ }\CommentTok{# create a matrix}
\ErrorTok{>}\StringTok{ }\NormalTok{z =}\StringTok{ }\KeywordTok{cbind}\NormalTok{(x,y)}
\NormalTok{>}\StringTok{ }\NormalTok{z}
     \NormalTok{x y}
\NormalTok{[}\DecValTok{1}\NormalTok{,] }\DecValTok{1} \DecValTok{2}
\NormalTok{[}\DecValTok{2}\NormalTok{,] }\DecValTok{2} \DecValTok{3}
\NormalTok{[}\DecValTok{3}\NormalTok{,] }\DecValTok{3} \DecValTok{4}
\NormalTok{>}\StringTok{ }\CommentTok{# create a matrix}
\ErrorTok{>}\StringTok{ }\NormalTok{w =}\StringTok{ }\KeywordTok{rbind}\NormalTok{(x,y)}
\NormalTok{>}\StringTok{ }\NormalTok{w}
  \NormalTok{[,}\DecValTok{1}\NormalTok{] [,}\DecValTok{2}\NormalTok{] [,}\DecValTok{3}\NormalTok{]}
\NormalTok{x    }\DecValTok{1}    \DecValTok{2}    \DecValTok{3}
\NormalTok{y    }\DecValTok{2}    \DecValTok{3}    \DecValTok{4}
\NormalTok{>}\StringTok{ }\CommentTok{# make w into data.frame}
\ErrorTok{>}\StringTok{ }\NormalTok{v =}\StringTok{ }\KeywordTok{as.data.frame}\NormalTok{(w)}
\NormalTok{>}\StringTok{ }\NormalTok{v}
  \NormalTok{V1 V2 V3}
\NormalTok{x  }\DecValTok{1}  \DecValTok{2}  \DecValTok{3}
\NormalTok{y  }\DecValTok{2}  \DecValTok{3}  \DecValTok{4}
\NormalTok{>}\StringTok{ }\CommentTok{# add a variable to v}
\ErrorTok{>}\StringTok{ }\NormalTok{v$New =}\StringTok{ }\KeywordTok{c}\NormalTok{(}\DecValTok{10}\NormalTok{,}\DecValTok{20}\NormalTok{)}
\NormalTok{>}\StringTok{ }\NormalTok{v}
  \NormalTok{V1 V2 V3 New}
\NormalTok{x  }\DecValTok{1}  \DecValTok{2}  \DecValTok{3}  \DecValTok{10}
\NormalTok{y  }\DecValTok{2}  \DecValTok{3}  \DecValTok{4}  \DecValTok{20}
\end{Highlighting}
\end{Shaded}

\subsection{Operations on
vectors/matrices}\label{operations-on-vectorsmatrices}

The operations \texttt{+}, \texttt{-}, \texttt{*} (multiplication),
\texttt{/}, and \texttt{\^{}} (power) operate \emph{entrwise} on vectors
and matrices.

\begin{Shaded}
\begin{Highlighting}[]
\NormalTok{>}\StringTok{ }\NormalTok{x =}\StringTok{ }\DecValTok{1}\NormalTok{:}\DecValTok{3}
\NormalTok{>}\StringTok{ }\NormalTok{x}
\NormalTok{[}\DecValTok{1}\NormalTok{] }\DecValTok{1} \DecValTok{2} \DecValTok{3}
\NormalTok{>}\StringTok{ }\NormalTok{y =}\StringTok{ }\DecValTok{2}\NormalTok{:}\DecValTok{4}
\NormalTok{>}\StringTok{ }\NormalTok{y}
\NormalTok{[}\DecValTok{1}\NormalTok{] }\DecValTok{2} \DecValTok{3} \DecValTok{4}
\NormalTok{>}\StringTok{ }\NormalTok{x+y}
\NormalTok{[}\DecValTok{1}\NormalTok{] }\DecValTok{3} \DecValTok{5} \DecValTok{7}
\NormalTok{>}\StringTok{ }\NormalTok{x-y}
\NormalTok{[}\DecValTok{1}\NormalTok{] -}\DecValTok{1} \NormalTok{-}\DecValTok{1} \NormalTok{-}\DecValTok{1}
\NormalTok{>}\StringTok{ }\NormalTok{x*y}
\NormalTok{[}\DecValTok{1}\NormalTok{]  }\DecValTok{2}  \DecValTok{6} \DecValTok{12}
\NormalTok{>}\StringTok{ }\NormalTok{x^}\DecValTok{2}
\NormalTok{[}\DecValTok{1}\NormalTok{] }\DecValTok{1} \DecValTok{4} \DecValTok{9}
\NormalTok{>}\StringTok{ }\KeywordTok{sqrt}\NormalTok{(x)}
\NormalTok{[}\DecValTok{1}\NormalTok{] }\FloatTok{1.000000} \FloatTok{1.414214} \FloatTok{1.732051}
\end{Highlighting}
\end{Shaded}

\subsection{Operations on
vectors/matrices}\label{operations-on-vectorsmatrices-1}

\begin{Shaded}
\begin{Highlighting}[]
\NormalTok{>}\StringTok{ }\NormalTok{x =}\StringTok{ }\DecValTok{1}\NormalTok{:}\DecValTok{3}
\NormalTok{>}\StringTok{ }\NormalTok{x}
\NormalTok{[}\DecValTok{1}\NormalTok{] }\DecValTok{1} \DecValTok{2} \DecValTok{3}
\NormalTok{>}\StringTok{ }\NormalTok{y =}\StringTok{ }\DecValTok{2}\NormalTok{:}\DecValTok{4}
\NormalTok{>}\StringTok{ }\NormalTok{y}
\NormalTok{[}\DecValTok{1}\NormalTok{] }\DecValTok{2} \DecValTok{3} \DecValTok{4}
\NormalTok{>}\StringTok{ }\CommentTok{# transpose of y}
\ErrorTok{>}\StringTok{ }\KeywordTok{t}\NormalTok{(y)}
     \NormalTok{[,}\DecValTok{1}\NormalTok{] [,}\DecValTok{2}\NormalTok{] [,}\DecValTok{3}\NormalTok{]}
\NormalTok{[}\DecValTok{1}\NormalTok{,]    }\DecValTok{2}    \DecValTok{3}    \DecValTok{4}
\NormalTok{>}\StringTok{ }\CommentTok{# matrix multiplication}
\ErrorTok{>}\StringTok{ }\NormalTok{x%*%y}
     \NormalTok{[,}\DecValTok{1}\NormalTok{]}
\NormalTok{[}\DecValTok{1}\NormalTok{,]   }\DecValTok{20}
\NormalTok{>}\StringTok{ }\NormalTok{x%*%}\KeywordTok{t}\NormalTok{(y)}
     \NormalTok{[,}\DecValTok{1}\NormalTok{] [,}\DecValTok{2}\NormalTok{] [,}\DecValTok{3}\NormalTok{]}
\NormalTok{[}\DecValTok{1}\NormalTok{,]    }\DecValTok{2}    \DecValTok{3}    \DecValTok{4}
\NormalTok{[}\DecValTok{2}\NormalTok{,]    }\DecValTok{4}    \DecValTok{6}    \DecValTok{8}
\NormalTok{[}\DecValTok{3}\NormalTok{,]    }\DecValTok{6}    \DecValTok{9}   \DecValTok{12}
\end{Highlighting}
\end{Shaded}

\begin{itemize}
\tightlist
\item
  \texttt{x\%*\%y}: sum of products of pairs of entries with matching
  indices.
\end{itemize}

\subsection{\texorpdfstring{\texttt{character} and
\texttt{factor}}{character and factor}}\label{character-and-factor}

\texttt{character} is often used to name variables or annotate figures,
and \texttt{factor} to describe values of a discrete random variable.

\begin{itemize}
\tightlist
\item
  A \texttt{character} can be created by enclosing a sequence of
  characters by quotes (\texttt{\textquotesingle{}\textquotesingle{}} or
  \texttt{"\ "}), or by concatenation (\texttt{c}) of characters
\item
  A \texttt{factor} can be created by the command \texttt{factor} or by
  coercing (\texttt{as.factor})
\end{itemize}

\subsection{\texorpdfstring{\texttt{character}}{character}}\label{character}

\begin{Shaded}
\begin{Highlighting}[]
\NormalTok{>}\StringTok{ }\NormalTok{a =}\StringTok{ 'This is'}
\NormalTok{>}\StringTok{ }\NormalTok{a}
\NormalTok{[}\DecValTok{1}\NormalTok{] }\StringTok{"This is"}
\NormalTok{>}\StringTok{ }\NormalTok{b =}\StringTok{ 'Stat 435'}
\NormalTok{>}\StringTok{ }\CommentTok{# both a and b are characters}
\ErrorTok{>}\StringTok{ }\CommentTok{# concatenate strings via paste}
\ErrorTok{>}\StringTok{ }\NormalTok{w=}\KeywordTok{paste}\NormalTok{(a,b,}\DataTypeTok{sep=}\StringTok{" "}\NormalTok{)}
\NormalTok{>}\StringTok{ }\NormalTok{w}
\NormalTok{[}\DecValTok{1}\NormalTok{] }\StringTok{"This is Stat 435"}
\NormalTok{>}\StringTok{ }\CommentTok{# concatenate strings via c}
\ErrorTok{>}\StringTok{ }\NormalTok{z =}\StringTok{ }\KeywordTok{c}\NormalTok{(a,b)}
\NormalTok{>}\StringTok{ }\CommentTok{#z is a vector whose entries are characters}
\ErrorTok{>}\StringTok{ }\NormalTok{z[}\DecValTok{1}\NormalTok{]}
\NormalTok{[}\DecValTok{1}\NormalTok{] }\StringTok{"This is"}
\NormalTok{>}\StringTok{ }\NormalTok{z[}\DecValTok{2}\NormalTok{]}
\NormalTok{[}\DecValTok{1}\NormalTok{] }\StringTok{"Stat 435"}
\end{Highlighting}
\end{Shaded}

\subsection{\texorpdfstring{\texttt{factor}}{factor}}\label{factor}

\begin{Shaded}
\begin{Highlighting}[]
\NormalTok{>}\StringTok{ }\NormalTok{x =}\StringTok{ }\KeywordTok{c}\NormalTok{(}\KeywordTok{rep}\NormalTok{(}\StringTok{'F'}\NormalTok{,}\DecValTok{6}\NormalTok{),}\KeywordTok{rep}\NormalTok{(}\StringTok{'M'}\NormalTok{,}\DecValTok{6}\NormalTok{))}
\NormalTok{>}\StringTok{ }\NormalTok{x}
 \NormalTok{[}\DecValTok{1}\NormalTok{] }\StringTok{"F"} \StringTok{"F"} \StringTok{"F"} \StringTok{"F"} \StringTok{"F"} \StringTok{"F"} \StringTok{"M"} \StringTok{"M"} \StringTok{"M"} \StringTok{"M"} \StringTok{"M"} \StringTok{"M"}
\NormalTok{>}\StringTok{ }\CommentTok{# obtain object type}
\ErrorTok{>}\StringTok{ }\KeywordTok{class}\NormalTok{(x)}
\NormalTok{[}\DecValTok{1}\NormalTok{] }\StringTok{"character"}
\NormalTok{>}\StringTok{ }\CommentTok{# convert x into factor}
\ErrorTok{>}\StringTok{ }\NormalTok{y =}\StringTok{ }\KeywordTok{as.factor}\NormalTok{(x)}
\NormalTok{>}\StringTok{ }\NormalTok{y}
 \NormalTok{[}\DecValTok{1}\NormalTok{] F F F F F F M M M M M M}
\NormalTok{Levels:}\StringTok{ }\NormalTok{F M}
\NormalTok{>}\StringTok{ }\KeywordTok{class}\NormalTok{(y)}
\NormalTok{[}\DecValTok{1}\NormalTok{] }\StringTok{"factor"}
\NormalTok{>}\StringTok{ }\KeywordTok{levels}\NormalTok{(y)}
\NormalTok{[}\DecValTok{1}\NormalTok{] }\StringTok{"F"} \StringTok{"M"}
\end{Highlighting}
\end{Shaded}

\subsection{\texorpdfstring{\texttt{factor}}{factor}}\label{factor-1}

\begin{Shaded}
\begin{Highlighting}[]
\NormalTok{>}\StringTok{ }\NormalTok{x =}\StringTok{ }\KeywordTok{c}\NormalTok{(}\KeywordTok{rep}\NormalTok{(}\DecValTok{1}\NormalTok{,}\DecValTok{6}\NormalTok{),}\KeywordTok{rep}\NormalTok{(}\DecValTok{2}\NormalTok{,}\DecValTok{6}\NormalTok{))}
\NormalTok{>}\StringTok{ }\NormalTok{x}
 \NormalTok{[}\DecValTok{1}\NormalTok{] }\DecValTok{1} \DecValTok{1} \DecValTok{1} \DecValTok{1} \DecValTok{1} \DecValTok{1} \DecValTok{2} \DecValTok{2} \DecValTok{2} \DecValTok{2} \DecValTok{2} \DecValTok{2}
\NormalTok{>}\StringTok{ }\KeywordTok{class}\NormalTok{(x)}
\NormalTok{[}\DecValTok{1}\NormalTok{] }\StringTok{"numeric"}
\NormalTok{>}\StringTok{ }\CommentTok{# convert x into factor}
\ErrorTok{>}\StringTok{ }\NormalTok{y =}\StringTok{ }\KeywordTok{as.factor}\NormalTok{(x)}
\NormalTok{>}\StringTok{ }\NormalTok{y}
 \NormalTok{[}\DecValTok{1}\NormalTok{] }\DecValTok{1} \DecValTok{1} \DecValTok{1} \DecValTok{1} \DecValTok{1} \DecValTok{1} \DecValTok{2} \DecValTok{2} \DecValTok{2} \DecValTok{2} \DecValTok{2} \DecValTok{2}
\NormalTok{Levels:}\StringTok{ }\DecValTok{1} \DecValTok{2}
\NormalTok{>}\StringTok{ }\KeywordTok{class}\NormalTok{(y)}
\NormalTok{[}\DecValTok{1}\NormalTok{] }\StringTok{"factor"}
\NormalTok{>}\StringTok{ }\KeywordTok{levels}\NormalTok{(y)}
\NormalTok{[}\DecValTok{1}\NormalTok{] }\StringTok{"1"} \StringTok{"2"}
\end{Highlighting}
\end{Shaded}

\subsection{\texorpdfstring{\texttt{data.frame}}{data.frame}}\label{data.frame}

Each column of a \texttt{data.frame} can have its own object type,
whereas all columns of a \texttt{matrix} have to have the same object
type.

\begin{Shaded}
\begin{Highlighting}[]
\NormalTok{>}\StringTok{ }\NormalTok{x =}\StringTok{ }\KeywordTok{data.frame}\NormalTok{(}\StringTok{"Variable1"}\NormalTok{=}\KeywordTok{c}\NormalTok{(}\DecValTok{1}\NormalTok{,}\DecValTok{3}\NormalTok{,}\DecValTok{5}\NormalTok{),}
\NormalTok{+}\StringTok{                "Variable2"}\NormalTok{=}\KeywordTok{as.factor}\NormalTok{(}\KeywordTok{c}\NormalTok{(}\StringTok{'A'}\NormalTok{,}\StringTok{'B'}\NormalTok{,}\StringTok{'C'}\NormalTok{)))}
\NormalTok{>}\StringTok{ }\NormalTok{x}
  \NormalTok{Variable1 Variable2}
\DecValTok{1}         \DecValTok{1}         \NormalTok{A}
\DecValTok{2}         \DecValTok{3}         \NormalTok{B}
\DecValTok{3}         \DecValTok{5}         \NormalTok{C}
\NormalTok{>}\StringTok{ }\CommentTok{# display data structure for x}
\ErrorTok{>}\StringTok{ }\KeywordTok{str}\NormalTok{(x)}
\StringTok{'data.frame'}\NormalTok{:}\StringTok{   }\DecValTok{3} \NormalTok{obs. of  }\DecValTok{2} \NormalTok{variables:}
\StringTok{ }\ErrorTok{$}\StringTok{ }\NormalTok{Variable1:}\StringTok{ }\NormalTok{num  }\DecValTok{1} \DecValTok{3} \DecValTok{5}
 \NormalTok{$}\StringTok{ }\NormalTok{Variable2:}\StringTok{ }\NormalTok{Factor w/}\StringTok{ }\DecValTok{3} \NormalTok{levels }\StringTok{"A"}\NormalTok{,}\StringTok{"B"}\NormalTok{,}\StringTok{"C"}\NormalTok{:}\StringTok{ }\DecValTok{1} \DecValTok{2} \DecValTok{3}
\end{Highlighting}
\end{Shaded}

\section{Basic operations on data}\label{basic-operations-on-data}

\subsection{Import and export data}\label{import-and-export-data}

\begin{itemize}
\tightlist
\item
  Rstudio's ``Import Data'' functionality provides a GUI to import data
  from \emph{CSV}, \emph{Excel}, etc, and implements the functionality
  of \texttt{read.csv}
\item
  The command \texttt{read.table} (\texttt{write.data}) imports (saves)
  data from (into) \texttt{.txt} or \texttt{.data} files
\item
  \texttt{save} (\texttt{save.image}) saves an object (the workspace)
  into an object
\end{itemize}

\subsection{A data set}\label{a-data-set}

The data set \texttt{Auto.data} (or file \texttt{Auto.csv}) contains
\texttt{mpg} (miles per gallon) for cars of different numbers of
\texttt{cylinders}, engine \texttt{displacement}, \texttt{horsepower},
manufactures (\texttt{name}), etc.

\begin{verbatim}
  mpg cylinders displacement horsepower weight acceleration
1  18         8          307        130   3504         12.0
2  15         8          350        165   3693         11.5
3  18         8          318        150   3436         11.0
4  16         8          304        150   3433         12.0
5  17         8          302        140   3449         10.5
6  15         8          429        198   4341         10.0
  year origin                      name
1   70      1 chevrolet chevelle malibu
2   70      1         buick skylark 320
3   70      1        plymouth satellite
4   70      1             amc rebel sst
5   70      1               ford torino
6   70      1          ford galaxie 500
\end{verbatim}

\subsection{Import and export data}\label{import-and-export-data-1}

Data are often stored in a directory or on a website. So, we need to
tell importing commands where data are. By default, R imports (exports)
data from (to) \emph{current working directory}.

\begin{Shaded}
\begin{Highlighting}[]
\NormalTok{>}\StringTok{ }\CommentTok{# show current working directory}
\ErrorTok{>}\StringTok{ }\KeywordTok{getwd}\NormalTok{()}
\NormalTok{[}\DecValTok{1}\NormalTok{] }\StringTok{"D:/MathRepo/stat435"}
\NormalTok{>}\StringTok{ }\NormalTok{AD1 =}\StringTok{ }\KeywordTok{read.table}\NormalTok{(}\StringTok{"Auto.data"}\NormalTok{,}\DataTypeTok{header =} \NormalTok{T,}\DataTypeTok{na.strings =} \StringTok{"?"}\NormalTok{)}
\NormalTok{>}\StringTok{ }\KeywordTok{class}\NormalTok{(AD1)}
\NormalTok{[}\DecValTok{1}\NormalTok{] }\StringTok{"data.frame"}
\NormalTok{>}\StringTok{ }\KeywordTok{dim}\NormalTok{(AD1)}
\NormalTok{[}\DecValTok{1}\NormalTok{] }\DecValTok{397}   \DecValTok{9}
\end{Highlighting}
\end{Shaded}

\begin{itemize}
\tightlist
\item
  \texttt{header\ =\ T}: the first line of the file contains the
  variable names
\item
  \texttt{na.strings\ =\ "?"}: the string \texttt{?} indicates a missing
  element in the data matrix
\end{itemize}

\subsection{Import and export data}\label{import-and-export-data-2}

Import data from a csv file online:

\begin{Shaded}
\begin{Highlighting}[]
\NormalTok{>}\StringTok{ }\NormalTok{AD2 =}\StringTok{ }\KeywordTok{read.csv}\NormalTok{(}\StringTok{"http://faculty.marshall.usc.edu/gareth-james/ISL/Auto.csv"}\NormalTok{,}
\NormalTok{+}\StringTok{                  }\DataTypeTok{header =} \NormalTok{T,}\DataTypeTok{na.strings =} \StringTok{"?"}\NormalTok{)}
\NormalTok{>}\StringTok{ }\KeywordTok{class}\NormalTok{(AD2)}
\NormalTok{[}\DecValTok{1}\NormalTok{] }\StringTok{"data.frame"}
\NormalTok{>}\StringTok{ }\KeywordTok{dim}\NormalTok{(AD2)}
\NormalTok{[}\DecValTok{1}\NormalTok{] }\DecValTok{397}   \DecValTok{9}
\NormalTok{>}\StringTok{ }\NormalTok{AD3 =}\StringTok{ }\NormalTok{AD2[}\DecValTok{1}\NormalTok{:}\DecValTok{2}\NormalTok{,}\DecValTok{1}\NormalTok{:}\DecValTok{5}\NormalTok{]}
\NormalTok{>}\StringTok{ }\NormalTok{AD3}
  \NormalTok{mpg cylinders displacement horsepower weight}
\DecValTok{1}  \DecValTok{18}         \DecValTok{8}          \DecValTok{307}        \DecValTok{130}   \DecValTok{3504}
\DecValTok{2}  \DecValTok{15}         \DecValTok{8}          \DecValTok{350}        \DecValTok{165}   \DecValTok{3693}
\NormalTok{>}\StringTok{ }\CommentTok{# export AD3}
\ErrorTok{>}\StringTok{ }\KeywordTok{write.csv}\NormalTok{(AD3,}\DataTypeTok{file=}\StringTok{"myFile.csv"}\NormalTok{)}
\end{Highlighting}
\end{Shaded}

\subsection{Save and load objects}\label{save-and-load-objects}

By default, R saves (loads) data to (from) \emph{current working
directory}.

\begin{Shaded}
\begin{Highlighting}[]
\NormalTok{>}\StringTok{ }\NormalTok{AD2 =}\StringTok{ }\KeywordTok{read.csv}\NormalTok{(}\StringTok{"http://faculty.marshall.usc.edu/gareth-james/ISL/Auto.csv"}\NormalTok{,}
\NormalTok{+}\StringTok{                  }\DataTypeTok{header =} \NormalTok{T,}\DataTypeTok{na.strings =} \StringTok{"?"}\NormalTok{)}
\NormalTok{>}\StringTok{ }\KeywordTok{save}\NormalTok{(AD2,}\DataTypeTok{file=}\StringTok{"Partial.RData"}\NormalTok{)}
\NormalTok{>}\StringTok{ }\KeywordTok{save.image}\NormalTok{(}\StringTok{"Everything.RData"}\NormalTok{)}
\NormalTok{>}\StringTok{ }\KeywordTok{load}\NormalTok{(}\StringTok{"Everything.RData"}\NormalTok{)}
\end{Highlighting}
\end{Shaded}

\textbf{Caution:} Loading into the current workspace anything that
contains variables with identical names to those in the currently
workspace will overwrite these objects in the current workspace.

\subsection{Check or clean data}\label{check-or-clean-data}

\begin{Shaded}
\begin{Highlighting}[]
\NormalTok{>}\StringTok{ }\NormalTok{AD1 =}\StringTok{ }\KeywordTok{read.table}\NormalTok{(}\StringTok{"Auto.data"}\NormalTok{,}\DataTypeTok{header =} \NormalTok{T,}\DataTypeTok{na.strings =} \StringTok{"?"}\NormalTok{)}
\NormalTok{>}\StringTok{ }\KeywordTok{class}\NormalTok{(AD1)}
\NormalTok{[}\DecValTok{1}\NormalTok{] }\StringTok{"data.frame"}
\NormalTok{>}\StringTok{ }\KeywordTok{dim}\NormalTok{(AD1)}
\NormalTok{[}\DecValTok{1}\NormalTok{] }\DecValTok{397}   \DecValTok{9}
\NormalTok{>}\StringTok{ }\CommentTok{# show variable names in data}
\ErrorTok{>}\StringTok{ }\KeywordTok{colnames}\NormalTok{(AD1)}
\NormalTok{[}\DecValTok{1}\NormalTok{] }\StringTok{"mpg"}          \StringTok{"cylinders"}    \StringTok{"displacement"} \StringTok{"horsepower"}  
\NormalTok{[}\DecValTok{5}\NormalTok{] }\StringTok{"weight"}       \StringTok{"acceleration"} \StringTok{"year"}         \StringTok{"origin"}      
\NormalTok{[}\DecValTok{9}\NormalTok{] }\StringTok{"name"}        
\NormalTok{>}\StringTok{ }\CommentTok{# remove na's from data}
\ErrorTok{>}\StringTok{ }\NormalTok{AD1a =}\StringTok{ }\KeywordTok{na.omit}\NormalTok{(AD1)}
\NormalTok{>}\StringTok{ }\KeywordTok{dim}\NormalTok{(AD1a)}
\NormalTok{[}\DecValTok{1}\NormalTok{] }\DecValTok{392}   \DecValTok{9}
\end{Highlighting}
\end{Shaded}

\emph{Note:} by default, a row with a missing value will be removed
\texttt{na.omit}

\subsection{Remarks}\label{remarks}

\begin{itemize}
\tightlist
\item
  There are other objects in R such as \texttt{list} and
  \texttt{tibble}, which are even more versatile than
  \texttt{data.frame}
\item
  There are other R commands to import data
\item
  The R package \texttt{dplyr} is a powerful tool for data processing,
  and the course contains supplementary videos for interested students
\end{itemize}

\section{Basic graphics commands}\label{basic-graphics-commands}

\subsection{\texorpdfstring{The \texttt{plot}
function}{The plot function}}\label{the-plot-function}

The \texttt{plot} function is a basic graphics command that creates 2D
plots. Its basic syntax is:

\begin{verbatim}
plot(x, y, ...)
\end{verbatim}

\begin{itemize}
\tightlist
\item
  \texttt{x}: the x coordinates of points in the plot
\item
  \texttt{y}: the y coordinates of points in the plot
\item
  \texttt{...}: other arguments such as \texttt{type} (what type of plot
  should be drawn), \texttt{main} (an overall title for the plot),
  \texttt{xlab} (a title for the x axis) and \texttt{ylab}(a title for
  the y axis)
\end{itemize}

\subsection{\texorpdfstring{The \emph{Auto} data
set}{The Auto data set}}\label{the-auto-data-set}

The data set \texttt{Auto.data} (or file \texttt{Auto.csv}) contains
\texttt{mpg} (miles per gallon) for cars of different numbers of
\texttt{cylinders}, engine \texttt{displacement}, \texttt{horsepower},
manufactures (\texttt{name}), etc.

\begin{verbatim}
  mpg cylinders displacement horsepower weight acceleration
1  18         8          307        130   3504         12.0
2  15         8          350        165   3693         11.5
3  18         8          318        150   3436         11.0
4  16         8          304        150   3433         12.0
5  17         8          302        140   3449         10.5
6  15         8          429        198   4341         10.0
  year origin                      name
1   70      1 chevrolet chevelle malibu
2   70      1         buick skylark 320
3   70      1        plymouth satellite
4   70      1             amc rebel sst
5   70      1               ford torino
6   70      1          ford galaxie 500
\end{verbatim}

\subsection{Illustration}\label{illustration}

\begin{Shaded}
\begin{Highlighting}[]
\NormalTok{>}\StringTok{ }\KeywordTok{plot}\NormalTok{(AD1a$displacement,AD1a$mpg,}\DataTypeTok{xlab=}\StringTok{"Engine displacement"}\NormalTok{,}
\NormalTok{+}\StringTok{      }\DataTypeTok{ylab=}\StringTok{"Miles per gallon"}\NormalTok{)}
\end{Highlighting}
\end{Shaded}

\begin{center}\includegraphics{Stat435LabNotes1_notes_files/figure-latex/plot1-1} \end{center}

\subsection{Illustration}\label{illustration-1}

\begin{Shaded}
\begin{Highlighting}[]
\NormalTok{>}\StringTok{ }\KeywordTok{attach}\NormalTok{(AD1a)}
\NormalTok{>}\StringTok{ }\KeywordTok{plot}\NormalTok{(cylinders,mpg,}\DataTypeTok{xlab=}\StringTok{"Number of cylinders"}\NormalTok{,}
\NormalTok{+}\StringTok{      }\DataTypeTok{ylab=}\StringTok{"Miles per gallon"}\NormalTok{)}
\end{Highlighting}
\end{Shaded}

\begin{center}\includegraphics{Stat435LabNotes1_notes_files/figure-latex/plot2-1} \end{center}

\subsection{\texorpdfstring{The \texttt{pairs}
function}{The pairs function}}\label{the-pairs-function}

The \texttt{pairs} function creates a \emph{scatterplot matrix}, i.e., a
scatterplot for every pair of variables for a given data set.

It can also produce a scatterplot matrix for a subset of variables (by
using the \texttt{\textasciitilde{}} operator).

\subsection{Illustration}\label{illustration-2}

\begin{Shaded}
\begin{Highlighting}[]
\NormalTok{>}\StringTok{ }\KeywordTok{attach}\NormalTok{(AD1a)}
\NormalTok{>}\StringTok{ }\KeywordTok{pairs}\NormalTok{(~mpg+displacement+horsepower,AD1a)}
\end{Highlighting}
\end{Shaded}

\begin{center}\includegraphics{Stat435LabNotes1_notes_files/figure-latex/plot3-1} \end{center}

\subsection{\texorpdfstring{The \texttt{identify}
function}{The identify function}}\label{the-identify-function}

\texttt{identify} reads the position of the graphics pointer when the
(first) mouse button is pressed. It then searches the coordinates given
in \texttt{x} and \texttt{y} for the point closest to the pointer.

If this point is close enough to the pointer, its index (as the index of
the row that contains the corresponding coordinate measurements for this
point) will be returned as part of the value of the call.

\subsection{Illustration}\label{illustration-3}

\begin{verbatim}
attach(AD1a)
plot(horsepower,mpg)
identify(horsepower,mpg)
# click on a point on the plot
# Hit Esc key once you have selected the point.
[1] 389 
# 389 is the index (i.e., row) number of the point 
# the Instructor selected
\end{verbatim}

\subsection{\texorpdfstring{The \texttt{hist}
function}{The hist function}}\label{the-hist-function}

The \texttt{hist} function creates histogram for a variable (or for
several variables).

\subsection{Illustration}\label{illustration-4}

\begin{Shaded}
\begin{Highlighting}[]
\NormalTok{>}\StringTok{ }\KeywordTok{attach}\NormalTok{(AD1a)}
\NormalTok{>}\StringTok{ }\KeywordTok{hist}\NormalTok{(horsepower,}\DataTypeTok{breaks=}\DecValTok{20}\NormalTok{)}
\end{Highlighting}
\end{Shaded}

\begin{center}\includegraphics{Stat435LabNotes1_notes_files/figure-latex/unnamed-chunk-16-1} \end{center}

\subsection{Illustration}\label{illustration-5}

\begin{Shaded}
\begin{Highlighting}[]
\NormalTok{>}\StringTok{ }\KeywordTok{attach}\NormalTok{(AD1a)}
\NormalTok{>}\StringTok{ }\KeywordTok{hist}\NormalTok{(horsepower,}\DataTypeTok{breaks=}\DecValTok{50}\NormalTok{)}
\end{Highlighting}
\end{Shaded}

\begin{center}\includegraphics{Stat435LabNotes1_notes_files/figure-latex/unnamed-chunk-17-1} \end{center}

\subsection{Remark}\label{remark}

The course Stat 437 covers approximately 3 weeks of materials for data
visualization.

\subsection{License and session
Information}\label{license-and-session-information}

\href{http://math.wsu.edu/faculty/xchen/stat412/LICENSE.html}{License}

\begin{Shaded}
\begin{Highlighting}[]
\NormalTok{>}\StringTok{ }\KeywordTok{sessionInfo}\NormalTok{()}
\NormalTok{R version }\FloatTok{3.5.0} \NormalTok{(}\DecValTok{2018-04-23}\NormalTok{)}
\NormalTok{Platform:}\StringTok{ }\NormalTok{x86_64-w64-mingw32/}\KeywordTok{x64} \NormalTok{(}\DecValTok{64}\NormalTok{-bit)}
\NormalTok{Running under:}\StringTok{ }\NormalTok{Windows }\DecValTok{7} \KeywordTok{x64} \NormalTok{(build }\DecValTok{7601}\NormalTok{) Service Pack }\DecValTok{1}

\NormalTok{Matrix products:}\StringTok{ }\NormalTok{default}

\NormalTok{locale:}
\NormalTok{[}\DecValTok{1}\NormalTok{] LC_COLLATE=English_United States}\FloatTok{.1252} 
\NormalTok{[}\DecValTok{2}\NormalTok{] LC_CTYPE=English_United States}\FloatTok{.1252}   
\NormalTok{[}\DecValTok{3}\NormalTok{] LC_MONETARY=English_United States}\FloatTok{.1252}
\NormalTok{[}\DecValTok{4}\NormalTok{] LC_NUMERIC=C                          }
\NormalTok{[}\DecValTok{5}\NormalTok{] LC_TIME=English_United States}\FloatTok{.1252}    

\NormalTok{attached base packages:}
\NormalTok{[}\DecValTok{1}\NormalTok{] stats     graphics  grDevices utils     datasets  methods  }
\NormalTok{[}\DecValTok{7}\NormalTok{] base     }

\NormalTok{other attached packages:}
\NormalTok{[}\DecValTok{1}\NormalTok{] knitr_1}\FloatTok{.21}

\NormalTok{loaded via a }\KeywordTok{namespace} \NormalTok{(and not attached):}
\StringTok{ }\NormalTok{[}\DecValTok{1}\NormalTok{] compiler_3}\FloatTok{.5.0}   \NormalTok{magrittr_1}\FloatTok{.5}     \NormalTok{tools_3}\FloatTok{.5.0}     
 \NormalTok{[}\DecValTok{4}\NormalTok{] htmltools_0}\FloatTok{.3.6}  \NormalTok{yaml_2}\FloatTok{.2.0}       \NormalTok{Rcpp_1}\FloatTok{.0.0}      
 \NormalTok{[}\DecValTok{7}\NormalTok{] codetools_0}\FloatTok{.2}\DecValTok{-15} \NormalTok{stringi_1}\FloatTok{.2.4}    \NormalTok{rmarkdown_1}\FloatTok{.11}  
\NormalTok{[}\DecValTok{10}\NormalTok{] stringr_1}\FloatTok{.3.1}    \NormalTok{xfun_0}\FloatTok{.4}         \NormalTok{digest_0}\FloatTok{.6.18}   
\NormalTok{[}\DecValTok{13}\NormalTok{] evaluate_0}\FloatTok{.12}   
\end{Highlighting}
\end{Shaded}


\end{document}
